\secciongrande{Programación de scripts dónde y cómo colocar los scripts}

\resumen{
  Los scripts de JavaScript se insertan en un documento HTML utilizando la etiqueta \texttt{<script>}. Pueden ubicarse tanto en el encabezado (entre \texttt{<head>} y \texttt{</head>}) como en el cuerpo (entre \texttt{<body>} y \texttt{</body>}). La etiqueta de apertura debe incluir el atributo \texttt{type="text/javascript"} para identificar el contenido como código JavaScript.
}

\nota{Para el examen}{
¿Cuál es la función principal del atributo \texttt{type="text/javascript"} en la etiqueta \texttt{<script>}? Indica al navegador que el contenido es código JavaScript y no HTML.
}

\begin{codigo}
\cmd{<script type="text/javascript">}
// Código JavaScript aquí
\cmd{</script>}
\end{codigo}



\textbf{Ubicación de scripts}
\lista{ |
  \textbf{Encabezado:} Se coloca entre \texttt{<head>} y \texttt{</head>}. Suele usarse para scripts que deben ejecutarse antes de que el contenido visible se cargue completamente. | | | | |
  \textbf{Cuerpo:} Se coloca entre \texttt{<body>} y \texttt{</body>}. Es la ubicación más común para scripts que interactúan con elementos visibles de la página. }


\subtituloA{Ejecución básica de un script}

\begin{codigo}
documento.write("Hola Onliner!");
\end{codigo}



\nota{Para el examen}{
¿Qué indica el punto y coma (\texttt{;}) al final de una instrucción en JavaScript? Señala al intérprete del navegador que la línea de código ha finalizado.
}